\documentclass[11pt]{article}
\usepackage[a4paper,margin=23mm]{geometry}
\usepackage{amsmath,amssymb,amsthm,mathtools,bm,mathrsfs}
\usepackage{booktabs,longtable,array}
\usepackage{enumitem}
\usepackage{xcolor}
\usepackage{microtype}
\usepackage[hidelinks]{hyperref}

\setlength{\emergencystretch}{3em}
\newtheorem{theorem}{Theorem}[section]
\newtheorem{proposition}[theorem]{Proposition}
\newtheorem{lemma}[theorem]{Lemma}
\newtheorem{corollary}[theorem]{Corollary}
\theoremstyle{definition}
\newtheorem{definition}[theorem]{Definition}
\newtheorem{gate}[theorem]{Fail-closed gate}
\theoremstyle{remark}
\newtheorem{remark}[theorem]{Remark}

\newcommand{\ii}{\mathrm i}
\newcommand{\dd}{\mathrm d}
\newcommand{\CC}{\mathbb C}
\newcommand{\RR}{\mathbb R}
\newcommand{\ZZ}{\mathbb Z}
\newcommand{\CP}{\mathbb {CP}}
\newcommand{\RP}{\mathbb {RP}}
\newcommand{\cO}{\mathcal O}
\newcommand{\cL}{\mathcal L}
\newcommand{\cP}{\mathcal P}
\newcommand{\cH}{\mathcal H}
\newcommand{\cM}{\mathcal M}
\newcommand{\cQ}{\mathcal Q}
\newcommand{\Tr}{\operatorname{Tr}}
\newcommand{\Hol}{\operatorname{Hol}}
\newcommand{\Pf}{\operatorname{Pf}}
\newcommand{\Det}{\operatorname{Det}}
\newcommand{\indtwo}{\operatorname{ind}_{2}}
\newcommand{\ch}{\operatorname{ch}}
\newcommand{\status}[1]{\textcolor{blue!60!black}{\textsf{#1}}}
\newcolumntype{L}[1]{>{\raggedright\arraybackslash}p{#1}}

\title{AP-E7 on the Actual \(SO(3)\) Soliton Orbit:\\
FR Torsion, Pfaffian Holonomy, CPT, and a Minimal\\
Rank-Three Resolution of the \(c_1=x+2y\) Obstruction}
\author{Route F research note}
\date{17 July 2026}

\begin{document}
\maketitle

\begin{abstract}
The canonical AP-E6 relaxed charged two-colour \(B=1\) profile has localized
collective orbit \(SO(3)=SU(2)/\ZZ_2\), not \(\CP^1\).  We therefore replace
the previously assumed free \(\cO(2)\) line by the exact classification of
lines on the actual orbit.  The cellular complex of \(\RP^3\) gives
\[
 H^2(SO(3);\ZZ)=\ZZ_2,\qquad
 t=\beta(a),\quad a\in H^1(SO(3);\ZZ_2),
\]
and the unique nontrivial flat line is
\(\cL_{\rm FR}=SU(2)\mathbin{\times}_{\rm sign}\CC\).  Its generator-loop
holonomy is \(-1\), its square is trivial, and its real form is the candidate
FR Pfaffian line.  A determinant obtained by squaring that real Pfaffian has
holonomy \(+1\).  However, the usual QCD WZW/FR rule gives
\((-1)^{N_cB}=+1\) for \(N_c=2,B=1\), conditionally selecting the trivial
bosonic sector.  The actual Yukawa mapping-torus mod-two index has not been
computed with a common microscopic regulator, so this standard selection is
not promoted to a same-operator Pfaffian theorem.

The torsion line is the correct alternative carrier on \(SO(3)\), but it is
not equivalent to \(\cO(2)\): its de Rham curvature vanishes and \(2t=0\),
whereas \(c_1(\cO(2))\) is a nonzero free class.  Their collective spectra
also disagree.  CPT inversion preserves both torsion classes.  We construct
one deck-compatible topological Real lift, while keeping the physical CPT
operator, its square, spectral cut, and Pfaffian orientation open.

In parallel we solve the earlier two-band obstruction at the level of bundle
mathematics.  On \(X=\CP^1_z\times\CP^1_w\), the three sections
\[
 s=(z_0w_0^2,\ z_1w_1^2,\ z_0w_1^2+z_1w_0^2)
\]
of \(\cO(1,2)\) have no common zero.  With
\(P_+=\overline{s}\,\overline{s}^{\dagger}/\lVert s\rVert^2\),
\[
 \cM=m(2P_+-\bm1_3),\qquad
 \operatorname{spec}\cM=\{+m,-m,-m\},\qquad
 c_1(E_+)=x+2y .
\]
The negative rank-two complement has
\(c_1(E_-)=-x-2y\) and \(c_2(E_-)=4xy\), so
\(E_+\oplus E_-=X\times\CC^3\) exactly.  This proves that rank three is
enough, and is minimal relative to the two-band obstruction.  It remains a
spectator: the \(\CP^1\) factor is not a normalizable modulus of the canonical
soliton and no charged two-colour Yukawa embedding has produced these three
bands.  Hence every physical closure and portal gate remains false.
\end{abstract}

\tableofcontents

\section{Scope, provenance, and theorem-level answer}

This note addresses two questions left by AP-E6.
\begin{enumerate}[label=(\roman*)]
\item What determinant or Pfaffian line can exist on the actual localized
  moduli \(SO(3)\), and what does CPT do to it?
\item Is the obstruction \(c_1^2\ne0\) intrinsic, or only a defect of the
  two-band asymptotic mass?
\end{enumerate}
The first question is answered exactly at the topological level and
fail-closed at the microscopic-regulator level.  The second has an explicit
rank-three success, fail-closed as a physical same-soliton construction.

The verifier imports the public AP-E6 Workline-B solver, calls
\begin{center}
\small
\nolinkurl{solve_relaxed_profile(box=16.0,sample_points=4097)},
\end{center}
and hard-asserts equality among its recomputed little-endian float64
\((r,F,s)\) digest and the Workline-B and Workline-C machine cards:
\begin{center}
\small
\nolinkurl{81104f59a5f3f4337739fc2a217cafc8da91581c9f1fb40b4fdba6ce0f948d59}.
\end{center}
The bound solution has
\(B=0.999999999686391\) and \(s(0)=0.917128964607058\).  Thus the topology
below is not attached to a separately solved or idealized hedgehog.

\begin{theorem}[Strongest defensible AP-E7 result]
Let \(\cQ=SO(3)\) be the localized orbit of the canonical AP-E6 \(B=1\)
profile.  Then:
\begin{enumerate}[label=(\alph*)]
\item complex line bundles on \(\cQ\) form \(\ZZ_2\); the nontrivial class is
  the flat FR sign line and has generator holonomy \(-1\);
\item its real sign form can be a Pfaffian line, but the actual family chooses
  it only if its mapping-torus mod-two index is odd;
\item inversion preserves the class and admits a topological Real lift, but
  topology does not determine the microscopic CPT regulator;
\item no line on \(\cQ\) is equivalent to the free \(\cO(2)\) line on
  \(\CP^1\);
\item a globally gapped rank-three Hermitian mass on the spectator
  \(\CP^1\times\CP^1\) realizes \(c_1(E_+)=x+2y\) exactly.
\end{enumerate}
Nevertheless,
\[
\begin{gathered}
\texttt{actual\_pfaffian\_torsion\_class\_selected=false},\\
\texttt{microscopic\_cpt\_regulator\_constructed=false},\\
\texttt{physical\_rank\_three\_yukawa\_embedding\_derived=false},\\
\texttt{lane\_closed=false},\qquad
\texttt{portal\_start\_allowed=false}.
\end{gathered}
\]
\end{theorem}

The determinant-line and eta-holonomy distinction follows the families
framework of Bismut--Freed and Dai--Freed
\cite{BismutFreed1986APE7,DaiFreed1994APE7}.  FR signs are configuration-space
quantization data rather than a consequence of merely writing a soliton
ansatz~\cite{FinkelsteinRubinstein1968APE7,Krusch2003APE7}.

\section{Exact topology of the physical orbit}

\subsection{Cellular computation of \texorpdfstring{\(H^2(SO(3);\ZZ)\)}{H2(SO(3);Z)}}

Use \(SO(3)\simeq\RP^3\), with one cell in each dimension \(0,1,2,3\).
For the usual cellular orientations the integral boundary is
\[
 d_k=1+(-1)^k,
 \qquad
 0\longrightarrow \ZZ\xrightarrow{0}\ZZ
 \xrightarrow{2}\ZZ\xrightarrow{0}\ZZ\longrightarrow0.
\]
Dualizing gives
\[
 0\longrightarrow \ZZ\xrightarrow{0}\ZZ
 \xrightarrow{2}\ZZ\xrightarrow{0}\ZZ\longrightarrow0,
\]
where the displayed terms now have cohomological degrees \(0,1,2,3\).
Consequently
\begin{equation}
 H^0=\ZZ,\qquad H^1=0,\qquad H^2=\ZZ/2,\qquad H^3=\ZZ.
 \label{eq:rp3-cohomology}
\end{equation}
In particular \(H^2_{\rm dR}(SO(3))=0\).  There is no continuous magnetic
flux or free Chern-number parameter on this orbit.

\subsection{The Bockstein and the unique nontrivial line}

Let \(a\) generate \(H^1(\RP^3;\ZZ_2)=\ZZ_2\), and use the coefficient
sequence
\[
 0\longrightarrow\ZZ\xrightarrow{\times2}\ZZ
 \longrightarrow\ZZ_2\longrightarrow0.
\]
Represent \(a\) by the mod-two cellular cochain \(1\), lift it to the integer
cochain \(A=1\), and note that \(\delta A=2\).  Dividing by two before taking
the integral class gives
\begin{equation}
 \beta(a)=\left[\frac{\delta A}{2}\right]=[1]=t,
 \qquad H^2(\RP^3;\ZZ)=\{0,t\},\qquad 2t=0.
 \label{eq:bockstein}
\end{equation}
Since complex line bundles are classified by \(H^2(-;\ZZ)\), these are the
only two topological possibilities.

The universal cover is
\[
 p:SU(2)=S^3\longrightarrow SO(3)=\RP^3,
 \qquad A\sim-A.
\]
The nontrivial line has the concrete associated-bundle model
\begin{equation}
 \cL_{\rm FR}=SU(2)\mathbin{\times}_{\ZZ_2,\mathrm{sign}}\CC,
 \qquad (A,z)\sim(-A,-z),
 \label{eq:fr-line}
\end{equation}
so a section lifted to \(SU(2)\) obeys \(\psi(-A)=-\psi(A)\).  Its pullback
to \(S^3\) is trivial because \(H^2(S^3;\ZZ)=0\), but its descent data are
not.  The real sign line
\[
 \cP_{\rm FR}=SU(2)\mathbin{\times}_{\mathrm{sign}}\RR
\]
has \(w_1(\cP_{\rm FR})=a\), and
\(\cP_{\rm FR}\otimes_{\RR}\CC=\cL_{\rm FR}\) has
\(c_1=\beta(a)=t\).

\subsection{Generator holonomy and rotor sectors}

A generator of \(\pi_1(SO(3))=\ZZ_2\) lifts to
\begin{equation}
 A(\tau)=
 \begin{pmatrix}e^{\ii\pi\tau}&0\\0&e^{-\ii\pi\tau}\end{pmatrix},
 \qquad 0\leq\tau\leq1,\qquad A(1)=-\bm1_2.
 \label{eq:generator-loop}
\end{equation}
Closing the lift uses the sign transition in~\eqref{eq:fr-line}; hence
\begin{equation}
 \Hol_\gamma(\cL_{\rm FR})=-1,
 \qquad \Hol_\gamma(\cL_{\rm FR}^{\otimes2})=+1.
 \label{eq:fr-holonomy}
\end{equation}
Peter--Weyl functions in spin \(j\) transform under the centre as
\((-1)^{2j}\).  The trivial sector therefore contains integer \(j\), while
the FR sign sector contains half-integer \(j\).  The lowest full rigid-rotor
block in the odd sector is \(j=1/2\) and has dimension
\((2j+1)^2=4\), not three.  Additional hedgehog constraints may reduce the
physical spin/isospin labeling, but they cannot turn the torsion line into
the \(\CP^1\) holomorphic space \(H^0(\CP^1,\cO(2))\).

\section{Determinant, Pfaffian, and the missing mapping-torus calculation}

\subsection{What topology determines}

For a real skew-adjoint Fredholm family with Pfaffian line \(\Pf\to\cQ\),
the loop sign is controlled by mod-two spectral flow, equivalently by the
appropriate mapping-torus mod-two index:
\begin{equation}
 \Hol_\gamma(\Pf)=(-1)^{\indtwo(\mathscr D_{Y_\gamma})}.
 \label{eq:pf-holonomy}
\end{equation}
On \(SO(3)\), equation~\eqref{eq:pf-holonomy} leaves exactly two options:
the trivial real line if the index is even, and \(\cP_{\rm FR}\) if it is
odd.  For the standard doubled skew matrix, the determinant line is the
Pfaffian square,
\begin{equation}
 \Det\simeq\Pf^{\otimes2},
 \qquad \Hol_\gamma(\Det)=+1
 \label{eq:det-pf-square}
\end{equation}
in either torsion sector.  This does not say that every complex chiral
determinant is automatically the square of one real Pfaffian; the real
structure and representation must first be specified.  A single
pseudoreal \(SU(2)\) Weyl doublet can carry the familiar mod-two anomaly,
whereas vectorlike doubling gives an even representation-parity control
\cite{Witten1982APE7}.

\subsection{Conditional two-colour selection}

The standard Skyrmion WZW/FR rotation rule gives
\begin{equation}
 \Hol_{2\pi}=(-1)^{N_cB}.
 \label{eq:ncb}
\end{equation}
For the present two-colour, unit-charge target,
\begin{equation}
 N_c=2,\qquad B=1,\qquad (-1)^{N_cB}=+1.
 \label{eq:two-colour-sign}
\end{equation}
Thus, \emph{if} the usual WZW derivation is realized by the same charged
two-colour regulator, it selects the trivial bosonic line, as expected for a
two-colour diquark baryon.  The AP-E6 declared constituent Dirac completion
also contains two Weyl doublets and has even mod-two representation parity.
These are strong consistency controls, not the requested operator theorem:
the common Euclidean/PV operator, its mapping-torus domain, eta prescription,
and Pfaffian orientation are still absent.  WZW quantization and global
current algebra supply the context for~\eqref{eq:ncb}~\cite{Witten1983APE7}.

\begin{gate}[Actual Pfaffian class]
The machine card sets
\begin{center}
\small
\nolinkurl{actual_yukawa_mapping_torus_mod2_index_computed=false}\\
\nolinkurl{actual_pfaffian_torsion_class_selected=false}.
\end{center}
The conditional value \(+1\) in~\eqref{eq:two-colour-sign} is not silently
substituted for a same-operator APS/PV computation.
\end{gate}

\section{CPT on the torsion line}

Use the base involution
\begin{equation}
 \iota:SO(3)\longrightarrow SO(3),\qquad \iota([A])=[A^{-1}].
 \label{eq:base-involution}
\end{equation}
Since inversion acts trivially on \(\pi_1(SO(3))=\ZZ_2\), it preserves both
flat characters.  Equivalently,
\[
 c_1(\overline{\cL}_{\rm FR})=-t=t,
\]
because \(2t=0\).  One explicit topological anti-linear lift is
\begin{equation}
 J[A,z]=[A^{-1},\overline z].
 \label{eq:real-lift}
\end{equation}
It is well-defined: replacing \((A,z)\) by \((-A,-z)\) sends the right-hand
side to \([-A^{-1},-\overline z]=[A^{-1},\overline z]\).  This particular
lift satisfies \(J^2=1\).

Existence is not uniqueness of the microscopic CPT operation.  A physical
antiunitary map must also act on spin, gauge, flavour and regulator spaces;
preserve or transform the operator domain and boundary conditions; choose a
spectral cut; and transport the determinant/Pfaffian orientation.  Its square
on the full Hilbert space can contain fermion parity or internal symmetries
not seen by the one-dimensional topological carrier.  We therefore record
\eqref{eq:real-lift} as a Real-line construction and keep
\(\texttt{microscopic\_cpt\_regulator\_constructed=false}\).

\section{Why torsion cannot replace \texorpdfstring{\(\cO(2)\)}{O(2)}}

Let \(u=c_1(\cO(1))\) generate \(H^2(\CP^1;\ZZ)=\ZZ\).  The old carrier has
\[
 c_1(\cO(2))=2u\ne0,
\]
a free class visible in de Rham cohomology and Berry curvature.  On the actual
orbit, every line has class \(0\) or \(t\), with
\[
 2t=0,
 \qquad t\otimes\RR=0,
 \qquad F_{\nabla_{\rm flat}}=0.
\]
There is no class-preserving identification \(2u\leftrightarrow t\).  More
physically, \(\cO(2)\) produces three Dolbeault sections on \(\CP^1\), while
the FR line imposes a centre parity on an \(SO(3)\) rotor.  Therefore the
torsion line is the correct \emph{classification replacement} on the true
moduli but not an equivalent realization of the AP-E4 three-state mechanism.

\section{Explicit minimal rank-three mass family}

\subsection{A basepoint-free map to \texorpdfstring{\(\CP^2\)}{CP2}}

Let
\[
 X=\CP^1_z\times\CP^1_w,
 \qquad x=c_1(\cO(1,0)),\qquad y=c_1(\cO(0,1)),
\]
so \(x^2=y^2=0\) and \(\int_Xxy=1\).  Consider
\begin{equation}
 \begin{aligned}
 s_0&=z_0w_0^2,\\
 s_1&=z_1w_1^2,\\
 s_2&=z_0w_1^2+z_1w_0^2.
 \end{aligned}
 \label{eq:sections}
\end{equation}
All three are sections of \(\cO(1,2)\).

\begin{lemma}[No common zero]
The sections in~\eqref{eq:sections} are basepoint-free.
\end{lemma}
\begin{proof}
If \(s_0=s_1=0\), then
\((z_0=0\ \text{or}\ w_0=0)\) and
\((z_1=0\ \text{or}\ w_1=0)\).  The branches
\(z_0=z_1=0\) and \(w_0=w_1=0\) are forbidden in projective space.  In the
remaining branch \(z_0=0,w_1=0\), one has
\(z_1w_0^2\ne0\), so \(s_2\ne0\).  In the branch
\(w_0=0,z_1=0\), one has \(z_0w_1^2\ne0\), again giving \(s_2\ne0\).
Thus no common zero exists.
\end{proof}

The normalized holomorphic column \(v=s/\lVert s\rVert\) defines a map
\(f:X\to\CP^2\).  Its span is the pulled-back tautological line
\[
 \operatorname{span}(v)=f^*\cO_{\CP^2}(-1),
 \qquad c_1(\operatorname{span}(v))=-(x+2y).
\]
This fixes the sign that a bare ``Veronese vector'' can easily obscure.
Complex conjugation reverses the Chern class.  Hence
\begin{equation}
 e_+=\frac{\overline{s}}{\lVert s\rVert},
 \qquad
 P_+=e_+e_+^\dagger,
 \qquad
 c_1(E_+)=x+2y.
 \label{eq:positive-projector}
\end{equation}

\subsection{Hermitian mass and uniform gap}

For any constant \(m>0\), define
\begin{equation}
 \cM(z,w)=m(2P_+(z,w)-\bm1_3).
 \label{eq:rank-three-mass}
\end{equation}
Because \(P_+^2=P_+=P_+^\dagger\) globally,
\begin{equation}
 \operatorname{spec}\cM=\{+m,-m,-m\},
 \qquad \cM^2=m^2\bm1_3.
 \label{eq:rank-three-spectrum}
\end{equation}
Thus the gap to zero is exactly \(m\) and the positive--negative band
separation is \(2m\), uniformly on all of \(X\).  Compactness plus the
basepoint-free lemma ensures the projector has no singular normalization.

\subsection{Complement bundle and K-theory identity}

Let \(E_-=E_+^\perp\), a rank-two bundle.  Since the ambient three-band
bundle is trivial,
\begin{equation}
 E_+\oplus E_-=X\times\CC^3,
 \qquad c(E_-)=(1+c_1(E_+))^{-1}.
\end{equation}
Writing \(\ell=x+2y\) and using \(\ell^2=4xy\),
\begin{equation}
 c_1(E_-)=-\ell=-(x+2y),
 \qquad c_2(E_-)=\ell^2=4xy.
 \label{eq:complement-chern}
\end{equation}
Equivalently in K-theory and rational cohomology
\begin{equation}
 [E_-]=3-[E_+],
 \qquad
 \ch(E_-)=3-e^\ell=2-\ell-\frac{\ell^2}{2}
 =2-(x+2y)-2xy .
 \label{eq:k-theory}
\end{equation}
These identities are the precise cancellation that a total trivial rank-three
bundle requires; no nonexistent ``complement line'' is introduced.  Standard
bundle classification and K-theory provide the general framework
\cite{Atiyah1989APE7,Hatcher2017APE7}.

\subsection{Minimality relative to the two-band no-go}

A gapped \(2\times2\) Hermitian mass with one positive and one negative band
defines a map to the Grassmannian
\(\mathrm{Gr}_1(\CC^2)=\CP^1\).  If \(u\in H^2(\CP^1;\ZZ)\), then
\(u^2=0\), so every positive eigenline pulled back from this target obeys
\[
 c_1(E_+)^2=0.
\]
The requested class obeys
\[
 (x+2y)^2=4xy\ne0.
\]
For three bands the target is \(\CP^2\), whose hyperplane class has nonzero
square.  Equations~\eqref{eq:positive-projector}--\eqref{eq:rank-three-spectrum}
therefore show both sufficiency and minimality: rank three is the first
nontrivial sign-split band number that can realize the class in a trivial
ambient bundle.

\section{Numerical and symbolic verification}

The deterministic verifier performs the following independent controls.
\begin{enumerate}[label=(\arabic*)]
\item It evaluates the \(\RP^3\) cellular boundaries \((0,2,0)\), the
  Bockstein lift \(\delta A/2=1\), and the FR character square.
\item It samples the path~\eqref{eq:generator-loop} at 65 points and checks
  unitarity, determinant one, endpoints, inversion, and deck compatibility.
\item It scans \(9\times12\) points on each \(\CP^1\), or 11,664 points on
  \(X\), at \(m=1.7\).  The sampled minimum is
  \(\lVert s\rVert^2=0.25\); projector, Hermiticity, and eigenvalue residuals
  are at floating-point roundoff, and the gap is \(1.7\).
\item Gauge-invariant triangulated Berry sums with 1,520 triangles per sphere
  give
  \[
  \frac{\ii}{2\pi}\int_{\CP^1_z}\Tr(P_+\dd P_+\wedge\dd P_+)=1,
  \qquad
  \frac{\ii}{2\pi}\int_{\CP^1_w}\Tr(P_+\dd P_+\wedge\dd P_+)=2
  \]
  to better than \(10^{-11}\).
\end{enumerate}
The matrix gap is an exact finite-band theorem.  It is not the Fredholm gap of
the spatial Yukawa--Callias operator.  General Callias theory requires the
actual asymptotic potential and operator class, not merely a target bundle
template~\cite{VanDenDungen2025APE7}.

\begin{table}[ht]
\centering
\caption{AP-E7 evidence boundary.}
\begin{tabular}{L{0.53\linewidth}L{0.30\linewidth}}
\toprule
Statement & Status\\
\midrule
\(H^2(SO(3);\ZZ)=\ZZ_2\), FR cover and both flat sectors & \status{exact}\\
FR sign-line and Pfaffian-candidate holonomy \(-1\) & \status{exact topology}\\
Standard \((-1)^{N_cB}=+1\) at \(N_c=2,B=1\) & \status{conditional UV rule}\\
Actual Yukawa mapping-torus mod-two index & \status{open}\\
Deck-compatible topological CPT Real lift & \status{constructed}\\
Microscopic CPT/PV regulator and Pfaffian orientation & \status{open}\\
Rank-three \(c_1(E_+)=x+2y\) mass and uniform matrix gap & \status{exact spectator}\\
Same-soliton rank-three Yukawa embedding and Fredholm gap & \status{open}\\
Degree-one Route-E portal & \status{forbidden}\\
\bottomrule
\end{tabular}
\end{table}

\section{Blockers and next discriminating calculations}

The result narrows the remaining work to operator-level questions.
\begin{enumerate}[label=\textbf{B\arabic*.},leftmargin=*]
\item Construct the common Euclidean/PV Yukawa operator along the nontrivial
  \(SO(3)\) loop and evaluate its mapping-torus \(\indtwo\) or eta phase.
  This decides between the two classified Pfaffian lines.
\item Define the regulator CPT action on its domain, spectral cut and
  orientation.  Then compare its transported Real structure with
  equation~\eqref{eq:real-lift}.
\item Derive three asymptotic fermion bands from an actual charged two-colour
  representation and match their normalized mass projector to
  equation~\eqref{eq:positive-projector}, including gauge covariance and
  heavy-threshold decoupling.
\item Replace the spectator \(\CP^1\) by a localized normalizable physical
  modulus, or formulate the family over the actual \(SO(3)\) orbit.  Without
  this step the rank-three success cannot close the same-soliton lane.
\item Prove a uniform Fredholm gap for that spatial operator.  Only then may
  the determinant/Pfaffian family be paired with the WZW descent and tested
  as a portal antecedent.
\end{enumerate}

\bibliographystyle{unsrt}
\bibliography{route_f/tex/ap_e7_so3_fr_rank3_family}

\end{document}
